% !TEX program = xelatex

% Le PDF produit doit être convertit en PNG. Pour l'instant, je le fais avec Gimp.

\documentclass[margin=1mm]{standalone}

\usepackage{fontspec}
\usepackage{xunicode}
\usepackage{xcolor}
\usepackage{tikz}
\usetikzlibrary{arrows.meta, positioning, decorations.pathreplacing,calligraphy}
\setmainfont[Ligatures=TeX]{Source Sans Pro}


\begin{document}

\begin{tikzpicture}[
  % Arrow style
  arrow/.style={->,>=stealth},
  % Boxes style
  box/.style={ rectangle, draw, rounded corners=2pt, minimum height=1cm, align=center, inner sep=6pt, text width=2cm, fill=white},
  txt/.style={ align=center, inner sep=6pt, text width=15mm},
]

  % Blocks
  \node[txt] (x) at ( 1,2) {Image $x$};
  \node[box] (d) at ( 4,2) {Degradation};
  \node[txt] (y) at ( 7,2) {Degraded image $y$};
  \node[box] (r) at (10,2) {Restoration};
  \node[txt] (z) at (13,2) {Restored image $\hat{x}$};

  % Arrows
  \draw [arrow] (x)--(d);
  \draw [arrow] (d)--(y);
  \draw [arrow] (y)--(r);
  \draw [arrow] (r)--(z);

  % Direct & inverse problems
  \draw [decorate, decoration={brace, mirror, amplitude=2mm}] (1.5,1.2)--(6.5,1.2);
  \node at ( 4,.5) {Direct problem};
  \draw [decorate, decoration={brace, mirror, amplitude=2mm}] (7.5,1.2)--(12.5,1.2);
  \node at (10,.5) {Inverse problem};

\end{tikzpicture}

\end{document}